\divider{Definitions}

Let \(\Omega\) denote either of the spatial domains introduced above.

\begin{definition}[Smoothness]\label{def:Smoothness}
A velocity\textendash pressure pair \((u,p)\) is smooth through a finite time \(T>0\) when
\[
  u\in C^\infty\!\left(\Omega\times[0,T];\mathbb R^3\right),
  \qquad
  p\in C^\infty\!\left(\Omega\times[0,T]\right).
\]
Equivalently, every mixed spatial and temporal derivative of finite order exists in the interior and extends continuously through \(T\).
The pair is smooth for all time when it is smooth through every finite \(T>0\).
\end{definition}

\begin{definition}[Continuation]\label{def:Continuation}
Suppose \((u,p)\) solves the Navier--Stokes equations on \(\Omega\times[0,T)\).
A continuation past \(T\) in a regularity class \(X\) is a pair \((\widetilde u,\widetilde p)\) that solves the same equations with the same viscosity on \(\Omega\times[0,T+\varepsilon)\) for some \(\varepsilon>0\), belongs to \(X\) there, and agrees with \((u,p)\) on \(\Omega\times[0,T)\).
The pressure normalization is held fixed when the two pairs are compared.
For the Clay problem, \(X=C^\infty\), so the continuation must remain smooth on the larger interval.
\end{definition}

\begin{definition}[Regularity]\label{def:Regularity}
Regularity is the general term for the function-space control available for a solution.
In the classical scale, \((u,p)\) has \(C^k\) regularity through \(T\) when every mixed spatial and temporal derivative of total order at most \(k\) exists in the interior and extends continuously through \(T\).
Thus \(C^0\) is continuity, and \(C^k\) requires continuous derivatives through total order \(k\).
Smoothness is the all-orders case, \(C^\infty\).
Sobolev and mixed space\textendash time bounds are other forms of regularity, measured by norms rather than by a pointwise derivative count.
\end{definition}

Regularity names the kind and degree of control available for the solution.
Smoothness selects the all-orders classical class \(C^\infty\), and continuation asks whether the Navier--Stokes solution can be extended past \(T\) while remaining in that class.
The Clay problem requires this past every finite \(T\).

At a finite endpoint \(T\), \(C^k\) regularity is lost when at least one required mixed derivative does not extend continuously to \(\Omega\times\{T\}\).
Near a finite spatial point, this may happen through unbounded growth or through bounded oscillation with no continuous limit.
A failure at any finite order prevents continuation in \(C^\infty\).
